\documentclass[11pt]{article}
\usepackage[a4paper,margin=1.1in]{geometry}
\usepackage[T1]{fontenc}
\usepackage[utf8]{inputenc}
\usepackage[ngerman]{babel}
\usepackage{lmodern}
\usepackage{amsmath,amssymb,amsthm,mathrsfs}
\usepackage{microtype}
\usepackage{fancyhdr}
\usepackage{enumitem}
\setlength{\parindent}{1.5em}
\setlength{\parskip}{0.18em}
\pagestyle{fancy}
\fancyhf{}
\lhead{Heinrich Weber, Lehrbuch der Algebra, Bd. II}
\rhead{Literal repair transcription}
\cfoot{\thepage}
\newcommand{\Om}{\Omega}
\newcommand{\Hm}{H_m}
\newcommand{\Norm}{\mathrm{N}}
\newcommand{\Normone}{\mathrm{N}_1}
\newcommand{\frp}{\mathfrak p}
\newcommand{\frP}{\mathfrak P}
\newcommand{\tg}{\operatorname{tg}}
\begin{document}

\section*{\S\,176. Der in $\Omega_m$ enthaltene reelle Körper $H_m$.}

Jede Zahl des Körpers $\Omega_m$ geht durch die Substitution $(r,r^{-1})$ in die conjugirt imaginäre Zahl über, die auch durch die Vertauschung $(i,-i)$ erhalten wird. Das Product zweier solcher conjugirt imaginärer Zahlen ist das Quadrat des absoluten Werthes einer jeden von ihnen.

Eine Zahl in $\Omega_m$ ist reell, wenn sie durch die Vertauschung $(r,r^{-1})$ ungeändert bleibt, und nur unter dieser Voraussetzung. Jede solche Zahl lässt sich rational durch die zweigliedrige Periode $r+r^{-1}$ ausdrücken, und gehört also einem reellen Körper vom Grade $\frac12\varphi(m)$ an, der ein Theiler von $\Omega_m$ ist. Diesen wollen wir den in $\Omega_m$ enthaltenen reellen Körper nennen und mit $H_m$ bezeichnen (obwohl auch noch andere reelle Körper, nämlich alle Theiler von $H_m$, in $\Omega_m$ enthalten sind).

Die $\frac12\varphi(m)=\frac12\mu$ Zahlen
\begin{equation}
1,\quad r+r^{-1},\quad r^2+r^{-2},\ldots,\quad r^{\frac12\mu-1}+r^{-\frac12\mu+1}
\tag{1}
\end{equation}
bilden eine Basis der ganzen Zahlen des Körpers $H_m$. Denn nach \S\,170, (7) ist
\begin{equation}
\begin{aligned}
\omega={}&x_0+x_1r+x_2r^2+\cdots+x_{\frac12\mu-1}r^{\frac12\mu-1}+x_{\frac12\mu}r^{\frac12\mu} \\
&+x_{-1}r^{-1}+x_{-2}r^{-2}+\cdots+x_{-\frac12\mu+1}r^{-\frac12\mu+1}
\end{aligned}
\tag{2}
\end{equation}
dann und nur dann eine ganze Zahl des Körpers $\Omega_m$, wenn die Coefficienten $x_0,x_1,\ldots$ ganze rationale Zahlen sind.

Die Bedingung dafür, dass diese Zahl dem Körper $H_m$ angehört, erhält man, wenn man die Substitution $(r,r^{-1})$ ausführt und die Differenz $=0$ setzt. Dividirt man noch durch $r-r^{-1}$, so ergiebt sich so:
\begin{align*}
0={}&(x_1-x_{-1})+(x_2-x_{-2})(r+r^{-1})+\cdots \\
&+(x_{\frac12\mu-1}-x_{-\frac12\mu+1})\frac{r^{\frac12\mu-1}-r^{-\frac12\mu+1}}{r-r^{-1}}
+x_{\frac12\mu}\frac{r^{\frac12\mu}-r^{-\frac12\mu}}{r-r^{-1}}.
\end{align*}
Die Divisionen lassen sich hier ausführen, und wenn man dann mit $r^{\frac12\mu-1}$ multiplicirt, so ergiebt sich für $r$ eine rationale Gleichung von niedrigerem als $\mu$tem Grade, die nur dann befriedigt sein kann, wenn alle ihre Coefficienten verschwinden. Dies führt zu den Gleichungen
\[
x_{\frac12\mu}=0,\qquad x_{\frac12\mu-1}=x_{-\frac12\mu+1},\ldots,\qquad x_1=x_{-1},
\]
und es ist also jede ganze Zahl $\omega$ des Körpers $H_m$ in der Form enthalten:
\[
\omega=x_0+x_1(r+r^{-1})+x_2(r^2+r^{-2})+\cdots+x_{\frac12\mu-1}(r^{\frac12\mu-1}+r^{-\frac12\mu+1}).
\]
Statt der Basis (1) kann man auch, wenn man
\[
\rho=r+r^{-1}
\]
setzt, die Potenzen von $\rho$:
\begin{equation}
1,\quad \rho,\quad \rho^2,\ldots,\quad \rho^{\frac12\mu-1}
\tag{3}
\end{equation}
als Basis der ganzen Zahlen von $H_m$ wählen. Denn die Grössen (3) lassen sich ganzzahlig durch die (1) ausdrücken und umgekehrt.

Die Zahl $\rho$ ist die Wurzel einer irreduciblen Gleichung vom Grade $\frac12\mu$, die wir mit $\psi(\rho)=0$ bezeichnen wollen. Ist $f(r)=0$ die Gleichung für $r$ (\S\,169), so besteht die identische Relation
\begin{equation}
f(x)=x^{\frac12\mu}\psi\left(x+\frac1x\right),
\tag{4}
\end{equation}
woraus durch Bildung der Ableitung
\begin{equation}
f'(r)=r^{\frac12\mu}\psi'(\rho)(1-r^{-2}).
\tag{5}
\end{equation}
Hieraus können wir die Grundzahl $\Delta_1$ des Körpers $H_m$ bilden, die, da $H_m$ nur reelle Zahlen enthält, positiv sein muss. Diese Grundzahl ist, wenn wir die Norm in Bezug auf $H_m$ mit $\Normone$ bezeichnen:
\[
\Delta_1=\pm\Normone\psi'(\rho).
\]
Ist aber $\Norm$ die in Bezug auf $\Omega_m$ gebildete Norm, so ist
\[
\Norm\psi'(\rho)=[\Normone\psi'(\rho)]^2=\Delta_1^2,
\]
und demnach ergiebt die Formel (5) mit Rücksicht auf \S\,169:
\[
\Delta=\Norm(1-r^{-2})\Delta_1^2.
\]
Nach \S\,169, (7) und (10) ist:
\begin{equation}
\begin{aligned}
\Norm(1-r^{-2})&=q &&\text{bei ungeradem }q,\\
&=4 &&\text{für }q=2,\\
\pm\Delta&=q\Delta_1^2 &&\text{bei ungeradem }q,\\
&=4\Delta_1^2 &&\text{für }q=2.
\end{aligned}
\tag{6}
\end{equation}
Setzt man darin
\[
\Delta=\pm q^{q^{\kappa-1}[\kappa(q-1)-1]},
\]
so folgt
\begin{equation}
\begin{aligned}
\Delta_1&=q^{\,q^{\kappa-1}\frac{q-1}{2}-\frac{q^{\kappa-1}+1}{2}} &&\text{bei ungeradem }q,\\
&=2^{(\kappa-1)2^{\kappa-2}-1} &&\text{für }q=2,
\end{aligned}
\tag{7}
\end{equation}
so dass $\Delta_1$ immer eine Potenz von $q$ ist.

\section*{\S\,177. Die Primideale im Körper $H_m$.}

Es ist jetzt die Frage zu untersuchen, in welcher Beziehung die Primideale des Körpers $H_m$ zu denen des Körpers $\Omega_m$ stehen.

Wenn $\frP$ irgend ein Primideal in $H_m$ vom Grade $f_1$ bedeutet, so muss $\frP$ Theiler einer natürlichen Primzahl $p$ sein, und $\frP$ kann also (nach \S\,171) nur dann durch die zweite oder eine höhere Potenz eines Primfactors $\frp$ in $\Omega_m$ theilbar sein, wenn $p=q$ ist.

Ist aber $\frP$ ein Theiler von $q$, so muss es eine Potenz von $\sigma=1-r$ sein. Wir setzen etwa
\[
\frP=\sigma^\lambda.
\]
Nehmen wir hiervon die Norm in Bezug auf $\Omega_m$, die das Quadrat der Norm in Bezug auf $H_m$ ist, so folgt (nach \S\,169)
\[
q^{2f_1}=q^\lambda,
\]
also $f_1=\frac12\lambda$. Da $f_1$ eine ganze Zahl ist, so muss $\lambda$ mindestens $=2$ sein. Es kann aber $\lambda$ auch nicht grösser als $2$ sein, da $\sigma^2$ mit der in $H_m$ enthaltenen ganzen Zahl $(1-r)(1-r^{-1})$ associirt ist, also $\sigma^2$ durch $\frP$ theilbar sein muss. Demnach ist $f_1=1$, und wir erhalten den Satz:

\begin{enumerate}[label=\arabic*.,leftmargin=2.2em]
\item Die Primzahl $q$ ist die $\frac12\mu$te Potenz einer in $H_m$ existirenden Primzahl ersten Grades.
\end{enumerate}

Es sei nun $p$ von $q$ verschieden, und $\frp$ ein Primfactor von $\frP$ im Körper $\Omega_m$ vom Grade $f$, der durch die Substitution $(r,r^{-1})$ in $\frp'$ übergeht. Dann ist $\frP$ sowohl durch $\frp$ als durch $\frp'$ theilbar, und andererseits ist $\frp\frp'$ in $H_m$ enthalten und also durch $\frP$ theilbar, und wir haben zu unterscheiden, ob $\frp,\frp'$ verschieden sind oder nicht.

Ist $\frp$ von $\frp'$ verschieden, so ist $\frP$ durch $\frp\frp'$ theilbar, und es folgt $\frP=\frp\frp'$, und durch Normbildung $f_1=f$.

Ist aber $\frp=\frp'$, so ist $\frP=\frp$, und die Normbildung ergiebt
\[
2f_1=f.
\]
Welcher dieser beiden Fälle eintritt, das hängt also davon ab, ob das Primideal $\frp$ die Substitution $(r,r^{-1})$ gestattet oder nicht, d. h. ob $(r,r^{-1})$ in der Gruppe, zu der das Primideal gehört, enthalten ist oder nicht. Nach \S\,172 ist dieser Unterschied dadurch bedingt, ob es irgend einen Exponenten $h$ giebt, für den die Congruenz
\[
p^h\equiv -1\pmod m
\]
erfüllt ist, oder ob es keinen solchen Exponenten giebt.

Wir fassen das Ergebniss folgendermaassen zusammen:

\begin{enumerate}[label=\arabic*.,leftmargin=2.2em,start=2]
\item Die von $q$ verschiedenen natürlichen Primzahlen $p$ zerfallen in zwei Arten $p_1,p_2$.

Die erste Art $p_1$ ist dadurch charakterisirt, dass für irgend einen Exponenten $h$
\[
p_1^h\equiv -1\pmod m,
\]
und diese Primzahlen zerfallen, wenn sie für den Modul $m$ zum Exponenten $f$ gehören und $\mu=ef$ gesetzt wird, im Körper $H_m$ in $e$ Primfactoren $\frac12 f$ten Grades. Ihre Primfactoren sind auch in $\Omega_m$ nicht weiter zerlegbar.

Die Primzahlen der zweiten Art $p_2$ sind dadurch bestimmt, dass keine ihrer Potenzen nach dem Modul $m$ mit $-1$ congruent wird, und sie zerfallen im Körper $H_m$ in $\frac12e$ Primfactoren $f$ten Grades. Ihre Primfactoren sind in $\Omega_m$ noch in je zwei Factoren zerlegbar.
\end{enumerate}

Bei den Primzahlen der ersten Art muss $f$ sicher eine gerade Zahl sein. Wenn $m$ ungerade ist, so genügt es auch, damit $p$ eine Primzahl von der ersten Art sei, dass sie zu einem geraden Exponenten gehöre; denn dann ist
\[
(p^{\frac12 f}-1)(p^{\frac12 f}+1)\equiv0\pmod m.
\]
Der erste dieser Factoren $p^{\frac12 f}-1$ ist aber nicht durch $m$ theilbar, weil sonst $p$ zum Exponenten $\frac12 f$ gehören würde. Folglich muss $p^{\frac12 f}+1$ durch $q$ theilbar sein. Hier können aber nicht beide Factoren $p^{\frac12 f}-1$ und $p^{\frac12 f}+1$ durch $q$ theilbar sein, weil sonst auch ihre Differenz, die $=2$ ist, durch $q$ theilbar wäre, und folglich ist $p^{\frac12 f}+1$ durch $m$ theilbar.

Ist aber $m$ eine Potenz von $2$, so ist die Congruenz
\[
p^h\equiv -1\pmod m
\]
nur dann möglich, wenn $p\equiv -1\pmod m$ ist. Denn jede gerade Potenz einer ungeraden Zahl ist $\equiv+1\pmod 8$ und kann also nicht $\equiv-1\pmod m$ sein. Ist aber $h$ ungerade, so ist
\[
\frac{p^h+1}{p+1}=p^{h-1}-p^{h-2}+\cdots+1
\]
selbst ungerade, und folglich ist $p^h+1$ durch keine höhere Potenz von $2$ theilbar, als $p+1$. Demnach können wir dem Satze 2. noch folgende nähere Bestimmung hinzufügen:

\begin{enumerate}[label=\arabic*.,leftmargin=2.2em,start=3]
\item Wenn $m$ ungerade ist, so gehören die Primzahlen $p_1$ zu einem geraden, die Primzahlen $p_2$ zu einem ungeraden Exponenten.

Ist $m$ gerade, so gehören die Primzahlen der Form $km-1$ zur ersten und alle anderen ungeraden Primzahlen zur zweiten Art.
\end{enumerate}

\section*{\S\,178. Die Einheiten des Körpers $H_m$.}

Die dem Körper $H_m$ angehörigen Einheiten sind von besonderer Wichtigkeit für die Anwendungen. Auf sie lassen sich auch die Einheiten des Körpers $\Omega_m$ zurückführen nach folgendem Satze:

\begin{enumerate}[label=\arabic*.,leftmargin=2.2em]
\item Jede Einheit des Körpers $\Omega_m$ ist das Product einer Einheitswurzel und einer reellen Einheit des Körpers $H_m$.
\end{enumerate}

Bezeichnen wir irgend eine Einheit des Körpers $\Omega_m$ mit $\mathscr E(r)$, so ist $\mathscr E(r^{-1})$ der conjugirt imaginäre Werth, der gleichfalls eine Einheit darstellt, und der Quotient $\mathscr E(r):\mathscr E(r^{-1})$ ist wieder eine Einheit, deren absoluter Werth
\begin{equation}
\sqrt{\frac{\mathscr E(r)}{\mathscr E(r^{-1})}\frac{\mathscr E(r^{-1})}{\mathscr E(r)}}
\tag{1}
\end{equation}
gleich $1$ ist, und folglich ist dieser Quotient eine Einheitswurzel und mithin $=\pm r^h$, worin $h$ irgend ein ganzzahliger Exponent ist (\S\,175, Satz 1., 2.). Wir haben also
\begin{equation}
\mathscr E(r)=\pm r^h\mathscr E(r^{-1}).
\tag{2}
\end{equation}

Es ist nun im weiteren Beweise ein kleiner Unterschied zu machen, je nachdem $m$ ungerade oder gerade ist.

Ist $m$ zunächst ungerade, so können wir in (2) den Exponenten $h$ gerade annehmen, da wir ihn nöthigenfalls durch $h+m$ ersetzen können. Ferner ist in der Formel (2) nur das obere Zeichen zulässig. Denn wenn das untere Zeichen stände, so würde folgen:
\[
\mathscr E(r)\equiv\mathscr E(1)\equiv-\mathscr E(1),\qquad 2\mathscr E(1)\equiv0\,[\mathrm{mod}\,(1-r)].
\]
Nach \S\,169 ist aber $1-r$ ein Theiler von $q$, also relativ prim zu $2$, und daher ist $\mathscr E(1)$ und folglich auch $\mathscr E(r)$ durch $(1-r)$ theilbar. Dies aber widerspricht der Voraussetzung, dass $\mathscr E(r)$ eine Einheit sein soll.

Demnach ergiebt sich aus (2):
\[
r^{-\frac12 h}\mathscr E(r)=r^{\frac12 h}\mathscr E(r^{-1}),
\]
woraus folgt, dass $r^{-\frac12 h}\mathscr E(r)$ eine reelle Einheit ist. Bezeichnen wir sie mit $e(r)$, so ist also
\begin{equation}
\mathscr E(r)=r^{\frac12 h}e(r),
\tag{3}
\end{equation}
woraus für diesen Fall das Theorem 1. bewiesen ist.

Wenn aber zweitens $m$ eine Potenz von $2$ ist, so ist $\mu=\frac12m$ und $r^\mu=-1$, und zugleich mit $r$ ist $-r$ eine $m$te Einheitswurzel. Wenn wir nöthigenfalls $h$ durch $h+\frac12m$ ersetzen, so können wir in (2) das obere Zeichen annehmen und setzen:
\begin{equation}
\mathscr E(r)=r^h\mathscr E(r^{-1}).
\tag{4}
\end{equation}
Es kommt jetzt darauf an, nachzuweisen, dass $h$ gerade sein muss.

Wenn wir $\mathscr E(r)$ durch die Basis $1,r,r^2,\ldots,r^{\mu-1}$ darstellen, so ergiebt sich
\begin{equation}
\mathscr E(r)=\sum_{s=0}^{\mu-1}x_sr^s,
\tag{5}
\end{equation}
worin die $x_s$ ganze rationale Zahlen sind. Nach (4) ist aber dann
\[
\sum_{s=0}^{\mu-1}x_sr^s=\sum_{s=1}^{\mu-1}x_sr^{h-s},
\]
und daraus folgt, dass $x_s=\pm x_{s'}$ ist, wenn $s+s'\equiv h\pmod\mu$. Ist nun $h$ ungerade, so ist aus zwei so verbundenen Zahlen die eine gerade, die andere ungerade (da $\mu$ eine Potenz von $2$ ist). Der Ausdruck (5) ergiebt also für $\mathscr E(r)$ ein Aggregat von Gliedern der Form
\[
x_s(r^s\pm r^{s'}).
\]
Es ist aber
\[
r^s\pm r^{s'}=r^s(1\pm r^{s'-s})=r^s(1\mp r^{\mu+s'-s})
\]
immer durch $1-r$ theilbar, und daraus würde folgen, dass $\mathscr E(r)$ durch $1-r$, was keine Einheit ist, theilbar wäre, was dem Begriffe der Einheit widerspricht. Also ist die Annahme unzulässig, dass in der Formel (4) der Exponent $h$ ungerade sei, und wir können wie oben
\[
\mathscr E(r)=r^{\frac12h}e(r)
\]
setzen, worin $e(r)$ eine reelle Einheit ist. Damit ist der Satz 1. bewiesen.

Wir wollen ein System von reellen Einheiten hervorheben: Nach \S\,169 sind die $\mu$ Grössen $1-r^n$ alle unter einander associirt. Demnach ist der Quotient
\[
\frac{1-r^n}{1-r^{n'}},
\]
worin $n,n'$ zwei relative Primzahlen zu $m$ bedeuten, eine Einheit. Daraus leitet man aber, wenn
\[
r=e^{\frac{2\pi i}{m}}
\]
gesetzt wird, die reelle Einheit
\begin{equation}
\frac{\sin\frac{n\pi}{m}}{\sin\frac{n'\pi}{m}}
=\frac{r^{\frac n2}-r^{-\frac n2}}{r^{\frac{n'}2}-r^{-\frac{n'}2}}
=r^{\frac{n'-n}{2}}\frac{1-r^n}{1-r^{n'}}
\tag{6}
\end{equation}
her. Ist $m$ gerade, so kann man $n'=n+\frac m2$ setzen, und erhält für diesen Fall die reellen Einheiten
\begin{equation}
\tau_n=\operatorname{tang}\frac{n\pi}{m},
\tag{7}
\end{equation}
worin $n$ jede ungerade Zahl bedeuten kann.

Ersetzt man $n$ durch $\frac12m-n$, so geht $\tau_n$ in den reciproken Werth über. Lässt man $n$ die Reihe der Zahlen
\[
1,3,5,\ldots,\frac m4-1
\]
durchlaufen, so erhält $\tau_n$ lauter positive echt gebrochene Werthe.

\end{document}
